problem:  
Let \( (M^{2n}, g, J, \omega) = (\mathbb{C}P^n, g_{\text{FS}}, J_{\text{std}}, \omega_{\text{FS}}) \) be the complex projective space equipped with the Fubini–Study metric of constant holomorphic sectional curvature \(c>0\).  
Let \(L^n \subset \mathbb{C}P^n\) be a compact, embedded, minimal Lagrangian submanifold, and consider its Hamiltonian second variation operator (Oh):
\[
\mathcal{L}_L(f) = \Delta f + n c\, f, \quad f \in C^\infty(L),
\]
where \(f\) represents an infinitesimal Hamiltonian variation preserving the Lagrangian condition.

Let \(\mathcal{K}_L = \ker(\mathcal{L}_L)\) denote the space of Hamiltonian Jacobi fields, and let \(\mathcal{K}_{\text{iso}}(L)\) denote the subspace generated by restrictions to \(L\) of ambient holomorphic Killing potentials (i.e., potentials associated to infinitesimal isometries of \(\mathbb{C}P^n\)).

**Conjectural Problem (Hamiltonian Rigidity Gap Conjecture):**  
Does there exist a constant \(C(n)>0\) such that for every compact minimal Lagrangian submanifold \(L^n \subset \mathbb{C}P^n\),
\[
\dim(\mathcal{K}_L / \mathcal{K}_{\text{iso}}(L)) \le C(n)\, b_1(L),
\]
and, moreover, if equality holds (i.e. the quotient kernel reaches this bound), must \(L\) be congruent to the standard totally geodesic real form \( \mathbb{R}P^n \subset \mathbb{C}P^n \)?

Equivalently: among all compact minimal Lagrangian immersions in complex projective space, does the space of *non-symmetric* Hamiltonian Jacobi fields admit a dimension bound controlled by the first Betti number, with rigidity at the symmetric model?

Why is it a "good" problem:  
1. **Elimination of trivial symmetries:**  
   By factoring out the contribution from ambient holomorphic Killing potentials, this conjecture isolates *intrinsic* Hamiltonian deformation directions. This refinement eliminates the ℝPⁿ counterexample issue and poses a meaningful stability–topology relation that could hold nontrivially.

2. **Analytic–Topological Bridge:**  
   The inequality proposes a deep connection between an *analytic object* (the Jacobi kernel representing infinitesimal minimal-Hamiltonian deformations) and a *topological invariant* (the first Betti number), potentially linking Hamiltonian stability theory to Hodge theoretic data on \(L\).

3. **Structural Motivation from Index Theory:**  
   In Kähler–Einstein manifolds, elliptic operators governing Hamiltonian stability have spectra influenced by curvature and harmonic forms. A bound of this type would conceptually parallel index bounds in minimal surface theory, but within the Lagrangian and symplectic context.

4. **Rigidity Mechanism:**  
   The equality case aims to identify when the only Hamiltonian Jacobi fields arise from ambient symmetries—precisely the situation for totally geodesic real forms—thus capturing a *Hamiltonian rigidity gap* phenomenon analogous to classical isometric rigidity theorems.

5. **Simplicity and Depth:**  
   The statement is self-contained, formally concise, and geometric, yet resolving it would require new synthesis among symplectic Hodge theory, spectral analysis, and the geometry of symmetric spaces. This “simple-to-state but deeply-entangled” nature exemplifies a high-quality mathematical problem.

Hence, this conjecture represents a plausible, rigorously stated, and geometrically motivated open problem sitting at the interface of spectral geometry, symplectic topology, and global differential geometry.